% GENERATED FILE. Edit canonical lessons, then run npm run generate:books.
% canonical-source-hash: fnv1a64:fe5162278891a95d
% canonical-lessons: HI-C384-tarakna, HI-C384-jamna, HI-C384-ubalnaa, HI-C384-galna, HI-C384-madhyam

\chapter{Heat and Cold}
\label{ch:hi-heat-and-cold}
\hlchaptermodality{384}

\begin{chapteropening}
\textbf{By the end of this chapter:} \emph{I can say what heat or cold will do to things.}

Complete the last lesson of chapter 384: I can say what heat or cold will do to things.
\end{chapteropening}

\section[\texorpdfstring{\hi{तड़कना} (taṛaknā)}{तड़कना (taṛaknā)}]{\hi{तड़कना} (taṛaknā) — to crack (with a sharp sound)}

\label{lesson:HI-C384-tarakna}



\begin{quote}
Before the new one: say the Hindi for to unravel, then the Hindi for to get wet.
\end{quote}

\subsection*{You'll want to know: \hi{तड़कना}}
\textbf{\hi{तड़कना}} — \emph{taṛaknā} — \textquotedblleft{}to crack (with a sharp sound)\textquotedblright{}.

\begin{quote}
\hi{अगर} \hi{ठंडे} \hi{गिलास} \hi{में} \hi{गरम} \hi{चाय} \hi{भरोगे} \hi{तो} \hi{गिलास} \hi{तड़क} \hi{जाएगा।}
\end{quote}

\begin{quote}
\emph{agar ṭhaṇḍe gilās meṁ garam cāy bharoge to gilās taṛak jāegā.}
\end{quote}

\begin{quote}
If you fill a cold glass with hot tea, the glass will crack.
\end{quote}

Past (it …): \hi{वह} \hi{गिलास} \hi{तड़क} \hi{गया}, \emph{vah gilās taṛak gayā}.

\begin{cousinweb}
\textbf{\hi{तड़कना}} (\emph{taṛaknā}) is a tadbhava, inherited from Sanskrit \emph{traṭ}, \textquotedblleft{}crackle\textquotedblright{} (CDIAL 5988).
\end{cousinweb}

\begin{grammarlens}[title={What you've built}]
One more word for saying what happens if.
\end{grammarlens}

\subsection*{Your turn}
\begin{itemize}
\raggedright
  \item \emph{Say it:} \emph{taṛaknā}
  \item \emph{Say it:} \emph{taṛaknā}, once more
  \item \emph{Say it:} \emph{agar ṭhaṇḍe gilās meṁ garam cāy bharoge to gilās taṛak jāegā.} — If you fill a cold glass with hot tea, the glass will crack.
  \item \emph{Recall:} say the Hindi for to unravel, then the Hindi for to get wet, then say \emph{taṛaknā} again
\end{itemize}

\begin{tcolorbox}[breakable,colback=teal!4,colframe=teal!35!black,title={Before you move on}]
What does \hi{तड़कना} mean? (\textquotedblleft{}To crack (with a sharp sound)\textquotedblright{}.) Say it once more.
\end{tcolorbox}
\section[\texorpdfstring{\hi{जमना} (jamnā)}{जमना (jamnā)}]{\hi{जमना} (jamnā) — to freeze, to set}

\label{lesson:HI-C384-jamna}



\begin{quote}
Before the new one: say the Hindi for to get wet, then the Hindi for to crack.
\end{quote}

\subsection*{You'll want to know: \hi{जमना}}
\textbf{\hi{जमना}} — \emph{jamnā} — \textquotedblleft{}to freeze, to set\textquotedblright{}.

\begin{quote}
\hi{अगर} \hi{रात} \hi{को} \hi{बहुत} \hi{ठंड} \hi{होगी} \hi{तो} \hi{पानी} \hi{जम} \hi{जाएगा।}
\end{quote}

\begin{quote}
\emph{agar rāt ko bahut ṭhaṇḍ hogī to pānī jam jāegā.}
\end{quote}

\begin{quote}
If it is very cold at night, the water will freeze.
\end{quote}

Past (it …): \hi{वह} \hi{पानी} \hi{जम} \hi{गया}, \emph{vah pānī jam gayā}.

\begin{cousinweb}
\textbf{\hi{जमना}} (\emph{jamnā}) is a tadbhava, inherited from Sanskrit \emph{yamyate}, \textquotedblleft{}is fixed\textquotedblright{} (CDIAL 10428).
\end{cousinweb}

\begin{grammarlens}[title={What you've built}]
One more word for saying what happens if.
\end{grammarlens}

\subsection*{Your turn}
\begin{itemize}
\raggedright
  \item \emph{Say it:} \emph{jamnā}
  \item \emph{Say it:} \emph{jamnā}, once more
  \item \emph{Say it:} \emph{agar rāt ko bahut ṭhaṇḍ hogī to pānī jam jāegā.} — If it is very cold at night, the water will freeze.
  \item \emph{Recall:} say the Hindi for to get wet, then the Hindi for to crack, then say \emph{jamnā} again
\end{itemize}

\begin{tcolorbox}[breakable,colback=teal!4,colframe=teal!35!black,title={Before you move on}]
What does \hi{जमना} mean? (\textquotedblleft{}To freeze, to set\textquotedblright{}.) Say it once more.
\end{tcolorbox}
\section[\texorpdfstring{\hi{उबलना} (ubalnā)}{उबलना (ubalnā)}]{\hi{उबलना} (ubalnā) — to come to the boil (of a liquid)}

\label{lesson:HI-C384-ubalnaa}



\begin{quote}
Before the new one: say the Hindi for to crack, then the Hindi for to freeze.
\end{quote}

\subsection*{You'll want to know: \hi{उबलना}}
\textbf{\hi{उबलना}} — \emph{ubalnā} — \textquotedblleft{}to come to the boil (of a liquid)\textquotedblright{}.

\begin{quote}
\hi{अगर} \hi{दूध} \hi{उबलेगा} \hi{तो} \hi{गैस} \hi{धीमी} \hi{कर} \hi{दो।}
\end{quote}

\begin{quote}
\emph{agar dūdh ublegā to gais dhīmī kar do.}
\end{quote}

\begin{quote}
If the milk boils, turn the gas down.
\end{quote}

Past (it …): \hi{वह} \hi{दूध} \hi{उबल} \hi{गया}, \emph{vah dūdh ubal gayā}.

\begin{cousinweb}
\textbf{\hi{उबलना}} (\emph{ubalnā}) is inherited from \emph{*ubbal}, \textquotedblleft{}boil up\textquotedblright{}, a form CDIAL reconstructs from its descendants (CDIAL 2339).
\end{cousinweb}

\begin{grammarlens}[title={What you've built}]
One more word for saying what happens if.
\end{grammarlens}

\subsection*{Your turn}
\begin{itemize}
\raggedright
  \item \emph{Say it:} \emph{ubalnā}
  \item \emph{Say it:} \emph{ubalnā}, once more
  \item \emph{Say it:} \emph{agar dūdh ublegā to gais dhīmī kar do.} — If the milk boils, turn the gas down.
  \item \emph{Recall:} say the Hindi for to crack, then the Hindi for to freeze, then say \emph{ubalnā} again
\end{itemize}

\begin{tcolorbox}[breakable,colback=teal!4,colframe=teal!35!black,title={Before you move on}]
What does \hi{उबलना} mean? (\textquotedblleft{}To come to the boil (of a liquid)\textquotedblright{}.) Say it once more.
\end{tcolorbox}
\section[\texorpdfstring{\hi{गलना} (galnā)}{गलना (galnā)}]{\hi{गलना} (galnā) — to soften (in cooking), to dissolve}

\label{lesson:HI-C384-galna}



\begin{quote}
Before the new one: say the Hindi for to freeze, then the Hindi for to come to the boil.
\end{quote}

\subsection*{You'll want to know: \hi{गलना}}
\textbf{\hi{गलना}} — \emph{galnā} — \textquotedblleft{}to soften (in cooking), to dissolve\textquotedblright{}.

\begin{quote}
\hi{अगर} \hi{दाल} \hi{नहीं} \hi{गलेगी} \hi{तो} \hi{उसे} \hi{थोड़ी} \hi{देर} \hi{और} \hi{पकाओ।}
\end{quote}

\begin{quote}
\emph{agar dāl nahīṃ galegī to use thoṛī der aur pakāo.}
\end{quote}

\begin{quote}
If the lentils don't soften, cook them a little longer.
\end{quote}

Past (it …): \hi{वह} \hi{दाल} \hi{गल} \hi{गई}, \emph{vah dāl gal gaī}.

\begin{cousinweb}
\textbf{\hi{गलना}} (\emph{galnā}) is a tadbhava, inherited from Sanskrit \emph{galati}, \textquotedblleft{}drips, oozes\textquotedblright{} (CDIAL 4074).
\end{cousinweb}

\begin{grammarlens}[title={What you've built}]
One more word for saying what happens if.
\end{grammarlens}

\subsection*{Your turn}
\begin{itemize}
\raggedright
  \item \emph{Say it:} \emph{galnā}
  \item \emph{Say it:} \emph{galnā}, once more
  \item \emph{Say it:} \emph{agar dāl nahīṃ galegī to use thoṛī der aur pakāo.} — If the lentils don't soften, cook them a little longer.
  \item \emph{Recall:} say the Hindi for to freeze, then the Hindi for to come to the boil, then say \emph{galnā} again
\end{itemize}

\begin{tcolorbox}[breakable,colback=teal!4,colframe=teal!35!black,title={Before you move on}]
What does \hi{गलना} mean? (\textquotedblleft{}To soften (in cooking), to dissolve\textquotedblright{}.) Say it once more.
\end{tcolorbox}
\section[\texorpdfstring{\hi{मध्यम} (madhyam)}{मध्यम (madhyam)}]{\hi{मध्यम} (madhyam) — medium, moderate}

\label{lesson:HI-C384-madhyam}



\begin{quote}
Before the new one: say the Hindi for to come to the boil, then the Hindi for to soften.
\end{quote}

\subsection*{You'll want to know: \hi{मध्यम}}
\textbf{\hi{मध्यम}} — \emph{madhyam} — \textquotedblleft{}medium, moderate\textquotedblright{}.

\begin{quote}
\hi{अगर} \hi{गैस} \hi{मध्यम} \hi{नहीं} \hi{होगी} \hi{तो} \hi{खाना} \hi{जल} \hi{जाएगा।}
\end{quote}

\begin{quote}
\emph{agar gais madhyam nahīṃ hogī to khānā jal jāegā.}
\end{quote}

\begin{quote}
If the gas isn't on medium, the food will burn.
\end{quote}

\begin{cousinweb}
\textbf{\hi{मध्यम}} (\emph{madhyam}) is a tatsama, borrowed from Sanskrit \emph{madhyama}, \textquotedblleft{}in the middle\textquotedblright{} (Monier-Williams).
\end{cousinweb}

\begin{grammarlens}[title={What you've built}]
One more word for saying what happens if.
\end{grammarlens}

\subsection*{Your turn}
\begin{itemize}
\raggedright
  \item \emph{Say it:} \emph{madhyam}
  \item \emph{Say it:} \emph{madhyam}, once more
  \item \emph{Say it:} \emph{agar gais madhyam nahīṃ hogī to khānā jal jāegā.} — If the gas isn't on medium, the food will burn.
  \item \emph{Recall:} say the Hindi for to come to the boil, then the Hindi for to soften, then say \emph{madhyam} again
\end{itemize}

\begin{tcolorbox}[breakable,colback=teal!4,colframe=teal!35!black,title={Before you move on}]
What does \hi{मध्यम} mean? (\textquotedblleft{}Medium, moderate\textquotedblright{}.) Say it once more.
\end{tcolorbox}
